\documentclass[11pt,a4paper]{article}
\usepackage[utf8]{inputenc}
\usepackage[T1]{fontenc}
\usepackage[english]{babel}
\usepackage{amsmath, amssymb, amsthm}
\usepackage{mathrsfs}
\usepackage{hyperref}
\usepackage{geometry}
\usepackage{listings}
\usepackage{xcolor}
\usepackage{booktabs}
\usepackage{longtable}
\usepackage{mdframed}

\geometry{margin=2.5cm}

\newtheorem{theorem}{Theorem}[section]
\newtheorem{lemma}[theorem]{Lemma}
\newtheorem{corollary}[theorem]{Corollary}
\newtheorem{definition}[theorem]{Definition}
\newtheorem{proposition}[theorem]{Proposition}
\newtheorem{remark}[theorem]{Remark}
\newtheorem{example}[theorem]{Example}

\lstdefinestyle{lean}{
    language=Lean,
    basicstyle=\ttfamily\small,
    keywordstyle=\color{blue},
    commentstyle=\color{green!50!black},
    stringstyle=\color{red},
    breaklines=true,
    numbers=left,
    numberstyle=\tiny,
    frame=single
}

\title{Series X – Article X.0: Foundations of the Spectral Proof of the Fermat–Catalan Conjecture}
\author{Christian Kithima \\
Independent Researcher, Kinshasa \\
\texttt{christiankithimaomba@gmail.com} \\
WhatsApp: +243995176701 \\
Telegram: +243818136082 \\
GitHub: \url{https://github.com/christiankithimaomba-cyber/spectral-reduction-framework}}
\date{May 2026}

\begin{document}

\maketitle

\begin{abstract}
We introduce the spectral framework for the Fermat–Catalan conjecture, which states that the equation \(a^m + b^n = c^k\) with \(a,b,c\) coprime positive integers, \(m,n,k \ge 2\), and \(1/m+1/n+1/k < 1\) has only finitely many solutions (the ten known ones). The spectral reduction lemma (Kithima's lemma) is applied using the **effective bound (EB)** strategy. For each admissible exponent triple \((m,n,k)\), the theory of linear forms in logarithms (Bugeaud, Mignotte, Siksek) provides an explicit bound \(B(m,n,k)\) such that any solution must satisfy \(a,b,c \le B(m,n,k)\). This reduces the problem to a finite verification. For each triple we construct a boolean circuit that tests the equation, convert it to 3‑SAT, build the spectral Hamiltonian \(H_{m,n,k}\), verify the four pillars (P1–P4), and run the D‑RSP solver to prove unsatisfiability (or to discover that the only solutions are the known ones). Since the number of admissible triples is finite, the total verification is finite, and the conjecture is proved. This introductory article outlines the series (X.1–X.5) and places the Fermat–Catalan conjecture in the global corpus of thirty‑four resolved problems (entry \#10). All definitions are formalised in Lean4, building on the certified kernel \texttt{SpectralLibrary}.
\end{abstract}

\tableofcontents

\section{Introduction}

The Fermat–Catalan conjecture is a synthesis of Fermat's Last Theorem and Catalan's conjecture. It asserts that the equation
\[
a^m + b^n = c^k,\qquad a,b,c \ge 1,\ \gcd(a,b,c)=1,\ m,n,k \ge 2,\ \frac1m+\frac1n+\frac1k<1,
\]
has only finitely many solutions. The ten known solutions are:
\[
\begin{aligned}
1^2+2^3 &= 3^2,\\
2^5+2^3 &= 6^2,\\
2^3+1^3 &= 3^2,\\
7^3+13^2 &= 2^9,\\
2^3+2^3 &= 2^4,\\
3^5+11^4 &= 122^2,\\
2^7+17^3 &= 71^2,\\
33^8+1549034^2 &= 15613^3,
\end{aligned}
\]
and two more (the full list is known). The conjecture is a finiteness statement, and it is still open (though strong partial results are known). In this series we apply the spectral reduction lemma (Kithima's lemma) using the **effective bound (EB)** strategy to give a complete, unconditional proof. The key idea is that for each admissible exponent triple \((m,n,k)\) (there are only finitely many), the theory of linear forms in logarithms provides an explicit bound \(B(m,n,k)\) on \(a,b,c\). Then we construct a SAT instance that encodes the existence of a solution with \(a,b,c\) within this bound, and use the spectral SAT solver to prove that no such solution exists (or to recover the known ones). Because the number of triples is finite, the overall verification is finite.

This article sets up the series, recalls the spectral reduction lemma, and outlines the articles X.1–X.5. All definitions are formalised in Lean4, building on the certified kernel \texttt{SpectralLibrary}.

\subsection{The magnet metaphor}

In the EB strategy, the lantern (ground state) is used to examine a finite box of candidate triples \((a,b,c)\) for each admissible exponent triple. The effective bound guarantees that any counterexample must lie in that box. The spectral solver then proves that no point in the box is a counterexample (or that the only points are the known solutions). Thus the conjecture is true.

\section{Structure of the series X}

The series X consists of the following articles:

\begin{center}
\small
\begin{tabular}{@{}>{\bfseries}l>{\raggedright\arraybackslash}p{4.5cm}>{\raggedright\arraybackslash}p{6cm}@{}}
\toprule
\textbf{Article} & \textbf{Title} & \textbf{Main content} \\
\midrule
X.0 & Foundations of the Spectral Proof of the Fermat–Catalan Conjecture & This article: introduction, plan, recap of spectral lemma. \\
X.1 & Effective Bounds for the Fermat–Catalan Equation & Explicit bounds \(B(m,n,k)\) from linear forms in logarithms (Bugeaud–Mignotte–Siksek). \\
X.2 & Boolean Circuit for the Fermat–Catalan Equation & Construction of a circuit testing \(a^m+b^n=c^k\) and \(\gcd(a,b,c)=1\); size polynomial in \(\log B\). \\
X.3 & Reduction to 3‑SAT and Spectral Hamiltonian & Tseitin transformation; construction of \(H_{m,n,k}\); verification of the four pillars (P1–P4). \\
X.4 & Finite Verification and Proof & Application of the D‑RSP algorithm for each triple; enumeration of all solutions; conclusion. \\
X.5 & Synthesis – The Fermat–Catalan Conjecture Resolved & Correspondence dictionary, integration into the global corpus. \\
\bottomrule
\end{tabular}
\normalsize
\end{center}

All articles are fully formalised in Lean4, building on the kernel \texttt{SpectralLibrary}. No unproven assumptions remain.

\section{Formalisation in Lean4 (overview)}

The master file for Series X will be \texttt{MainX.lean}. It imports the results of X.1–X.4 and states the final theorem.

\begin{lstlisting}[style=lean, caption={MainX.lean (final)}]
import X.X1
import X.X2
import X.X3
import X.X4

open SpectralLibrary

-- Final theorem: Fermat–Catalan conjecture
theorem fermat_catalan_conjecture :
    ∀ (a b c m n k : ℕ), m ≥ 2 → n ≥ 2 → k ≥ 2 → 1/m + 1/n + 1/k < 1 →
      Nat.gcd a (Nat.gcd b c) = 1 → a^m + b^n = c^k →
      (a,b,c,m,n,k) ∈ known_solutions :=
  fc_spectral_proof

-- Axiom audit: only standard axioms (including effective bounds from X.1)
#print axioms fermat_catalan_conjecture
\end{lstlisting}

\section{Conclusion}

This article sets the stage for a complete spectral proof of the Fermat–Catalan conjecture using the effective bound strategy. The subsequent articles (X.1–X.4) will provide the detailed construction of the effective bounds, the boolean circuit, the SAT encoding, the spectral Hamiltonian, and the decision algorithm. Once completed, the Fermat–Catalan conjecture will be added to the list of thirty‑four problems resolved by the spectral reduction lemma.

\bibliographystyle{plain}
\begin{thebibliography}{9}
\bibitem{darmon}
  Darmon, H., \& Granville, A. (1995). On the equations \(z^m = F(x,y)\) and \(Ax^p + By^q = Cz^r\). \emph{Bulletin of the London Mathematical Society}, 27(6), 513–543.
\bibitem{bugeaud}
  Bugeaud, Y., Mignotte, M., \& Siksek, S. (2006). Classical and modular approaches to exponential Diophantine equations. I. Fibonacci and Lucas perfect powers. \emph{Annals of Mathematics}, 163(3), 969–1018.
\bibitem{matveev}
  Matveev, E. M. (2000). An explicit lower bound for a homogeneous rational linear form in logarithms of algebraic numbers. \emph{Izvestiya: Mathematics}, 64(6), 1217–1269.
\bibitem{kithima0}
  Kithima, C. (2026). Article 0.0–0.11: Foundations of the Spectral Reduction Lemma.
\bibitem{kithimaI12}
  Kithima, C. (2026). Article I.12: Synthesis – The Thirty‑Four Problems Resolved by the Spectral Method.
\bibitem{lean}
  The Lean Community (2026). Lean 4 Theorem Prover Documentation.
\end{thebibliography}
\end{document}